% ============================================================
% Chapter V — Discussion
% ============================================================
\section{Discussion}

\subsection{Interpreting Results}

The multinomial logistic regression model identified two dominant predictors
of finishing tier: GridPosition and the Degradation $\times$ Fuel interaction
term. Together, these variables capture the two axes of F1 performance---where
you start (baseline pace) and how you manage the race (strategic execution).

The overwhelming significance of GridPosition ($p = 0.000$ for both tiers)
confirms that qualifying performance sets the foundation for race outcomes.
The coefficient of $-0.80$ for the Podium tier translates to a 55\% reduction
in podium odds per grid position---a finding consistent with the aerodynamic
and track-position dependency of modern F1 where overtaking is difficult at
most circuits. This result aligns with the broader motorsport analytics
literature showing that starting position is the single strongest predictor
of finishing order \parencite{kuhn2022f1}.

The significance of the Degradation $\times$ Fuel interaction term
($p = 0.0095$ for Podium) is the central finding of this study. Its positive
coefficient indicates that the combined physical state of high tire age and
high fuel mass---the opening laps of a race where the car is heaviest and
tires are newest---is associated with podium outcomes. This is
mechanistically coherent: top-tier teams extract competitive lap times even
under the most physically demanding conditions (heavy car, fresh but
un-optimized tires), while midfield competitors lose proportionally more
performance to the same stresses.

The individual significance of Fuel\_Mass\_Proxy for both tiers
($p = 0.0009$ Podium, $p = 0.0003$ Point Scorer) with negative coefficients
reflects the physical reality that lighter cars are faster. As fuel burns off
and the proxy decreases, the log-odds of being in a higher tier increase.
However, this variable is partially confounded with race progression (lap
number), making the interaction term a cleaner measure of the tire--fuel
coupling.

The non-significance of tire compound dummies (Tire\_SOFT, Tire\_MEDIUM for
Podium) suggests that \textit{which} compound a driver uses matters less than
\textit{how} the tire degrades over its stint. This supports the theoretical
position that degradation dynamics---not static compound labels---drive
performance differentiation.

\subsection{Relating to Literature}

The findings connect to several streams in the reviewed literature.

The Box's M test rejection of covariance homogeneity aligns with
\textcite{mendoza2023}, who similarly found that performance-tier groups
exhibit unequal covariance structures, necessitating multinomial logistic
regression over discriminant analysis. This pattern appears to be general
across classification problems where group means and variances co-vary.

The effectiveness of the interaction term supports the feature engineering
philosophy advocated by \textcite{santos2020}: domain-specific interaction
terms that represent real physical relationships outperform raw variables.
The Degradation $\times$ Fuel product is not an arbitrary statistical
construction---it represents a real force coupling that teams actively manage
during races.

The LOGO-CV validation approach follows the recommendations of
\textcite{varoquaux2019} and \textcite{wong2023}, who demonstrated that
group-aware validation produces more conservative and realistic performance
estimates. Our 62.36\% accuracy, while moderate, is an honest estimate that
accounts for the grouped structure of racing data.

The tire degradation findings are consistent with \textcite{baldi2022}, who
showed that degradation effects are nonlinear. Our linear tire age predictor
may understate late-stint effects, which could partly explain the lower
recall for the Point Scorer tier---midfield drivers in the latter half of
stints may exhibit degradation patterns that the linear model cannot capture.

\subsection{Study Limitations}

Several limitations should be acknowledged:

\begin{enumerate}
  \item \textbf{Fuel mass approximation.} The fuel proxy assumes a uniform
    linear burn-off from 110\,kg to 0\,kg across all teams and drivers.
    Actual fuel loads and consumption rates are proprietary and vary by team
    strategy. Teams that under-fuel their cars (running lighter than the
    proxy suggests) would appear slower in the model than they physically
    are, potentially biasing tier predictions.

  \item \textbf{Three-race scope.} While the three circuits cover distinct
    architectural types, three races represent only 13\% of the 2023
    calendar. Rain-affected races, sprint weekends, high-degradation circuits
    like Silverstone, and high-downforce tracks like Hungary are not
    represented. The model's generalizability to these conditions is unknown.

  \item \textbf{Unmodeled race dynamics.} DRS trains, safety car deployments,
    virtual safety car periods, undercut and overcut strategies, and team
    orders are not explicitly modeled. These factors create significant
    variance in the mid-field, which likely explains the low Point Scorer
    recall (0.45).

  \item \textbf{Linear tire age.} Tire degradation is known to follow a
    nonlinear trajectory with a performance ``cliff'' after a critical wear
    threshold \parencite{baldi2022}. The linear tire age variable cannot
    capture this nonlinearity, potentially understating the effect of
    late-stint degradation on finishing tier.

  \item \textbf{DNF protocol limitation.} Reassigning all DNF laps to ``No
    Points'' removes potentially informative telemetry from drivers who were
    running competitively before retiring. A more nuanced approach---such as
    using running position at the time of retirement---could improve model
    training.
\end{enumerate}

\subsection{Practical Impacts}

Despite its limitations, the model has practical relevance:

\begin{itemize}
  \item \textbf{Strategy validation:} The significance of the interaction term
    provides quantitative evidence that tire--fuel management is a
    statistically distinct driver of podium outcomes. Teams can use this to
    validate strategic simulations that model these forces jointly rather than
    independently.

  \item \textbf{Feature engineering template:} The mean-centered interaction
    approach is transferable to any domain where multiple physical processes
    co-occur and compound---battery degradation and temperature in EVs,
    material fatigue and load stress in structural engineering, or player
    fatigue and match intensity in team sports.

  \item \textbf{Modeling baseline for F1 analytics:} The 62.36\% LOGO-CV
    accuracy establishes a transparent, interpretable baseline against which
    more complex models (gradient-boosted trees, neural networks) can be
    benchmarked. Any black-box model should clear this bar while demonstrating
    that its additional complexity is justified.
\end{itemize}
